\Article{12}
Let \(x,x',x''\) etc. denote indefinitely the mutually independent errors of observations of the same kind, whose relative probabilities are expressed by the prefixed characteristic \(\varphi\); and let \(y\) denote a given rational function of the indeterminates \(x,x',x''\) etc.: the multiple integral (I)
\[
\int\varphi x.\varphi x'.\varphi x''\mathord{.}\,\mathord{.}\,\mathord{.}\,\mathord{.}\,\mathord{.}\,\mathrm{d}x.\mathrm{d}x'.\mathrm{d}x''\mathord{.}\,\mathord{.}\,\mathord{.}\,\mathord{.}
\]
extended over all values of the indeterminates \(x,x',x''\) for which the value of \(y\) falls within the given limits \(0\) and \(\eta\) will express the probability of a value of \(y\) lying indefinitely within \(0\) and \(\eta\). Evidently this integral will be a function of \(\eta\), whose differential we shall put \(=\psi\eta.\mathrm{d}\eta\), so that the integral itself becomes equal to the integral \(\int\psi\eta.\mathrm{d}\eta\) beginning at \(\eta=0\). In this way the characteristic \(\psi\eta\) must at the same time be regarded as expressing the relative probability of any value of \(y\). Since \(x\) can be regarded as a function of the indeterminates \(y,x',x''\) etc., which we shall put
\[
=f(y,x',x''\mathord{.}\,\mathord{.}\,\mathord{.}\,\mathord{.})
\]
the integral (I) will become
\[
\begin{aligned}
=\int\varphi.f(y,x',x''\mathord{.}\,\mathord{.}\,\mathord{.}).\frac{\mathrm{d}f(y,x',x''\mathord{.}\,\mathord{.}\,\mathord{.}\,\mathord{.})}{\mathrm{d}y}.{}\\
\varphi x'.\varphi x''\mathord{.}\,\mathord{.}\,\mathord{.}\,\mathord{.}\,\mathrm{d}y.\mathrm{d}x'.\mathrm{d}x''\mathord{.}\,\mathord{.}\,\mathord{.}\,\mathord{.}
\end{aligned}
\]
where \(y\) must be extended from \(y=0\) to \(y=\eta\), while the remaining indeterminates must range over all values to which a real value of \(f(y,x',x''\mathord{.}\,\mathord{.}\,\mathord{.}\,\mathord{.})\) corresponds. From this \SourceLeaf{018}it follows that
\[
\begin{aligned}
\psi y=\int\varphi.f(y,x',x''\mathord{.}\,\mathord{.}\,\mathord{.}\,\mathord{.}).\frac{\mathrm{d}f(y,x',x''\mathord{.}\,\mathord{.}\,\mathord{.}\,\mathord{.})}{\mathrm{d}y}.{}\\
\varphi x'.\varphi x''\mathord{.}\,\mathord{.}\,\mathord{.}\,\mathord{.}\,\mathrm{d}x'.\mathrm{d}x''\mathord{.}\,\mathord{.}\,\mathord{.}\,\mathord{.}
\end{aligned}
\]
the integration, in which \(y\) must be regarded as constant, extending over all values of the indeterminates \(x',x''\) etc. that give \(f(y,x',x''\mathord{.}\,\mathord{.}\,\mathord{.}\,\mathord{.})\) a real value.

\Article{13}
To carry out this integration in practice, knowledge of the function \(\varphi\), which is usually unknown, would be required: indeed, even if this function were known, in most cases the integration would exceed the powers of analysis. This being so, we shall be unable to assign the probability of the individual values of \(y\): but the case will be different if only the mean value of \(y\) is wanted, which arises from the integration of \(\int y\psi y.\mathrm{d}y\) extended over all values of \(y\) that it can attain. And since, evidently, for all values that \(y\) cannot attain, either because of the nature of the function it expresses (e.g. for negative values, if we had \(y=xx+x'x'+x''x''+\text{etc.}\)), or because certain limits are set to the errors \(x,x',x''\) etc. themselves, we must put \(\psi y=0\), the result will evidently be the same if that integration is extended over all real values of \(y\), namely from \(y=-\infty\) to \(y=+\infty\). Now the integral \(\int y\psi y.\mathrm{d}y\) taken between definite limits, namely from \(y=\eta\) to \(y=\eta'\), is equal to the integral
\[
\begin{aligned}
\int y\varphi.f(y,x',x''\mathord{.}\,\mathord{.}\,\mathord{.}\,\mathord{.})\frac{\mathrm{d}f(y,x',x''\mathord{.}\,\mathord{.}\,\mathord{.}\,\mathord{.})}{\mathrm{d}y}.{}\\
\varphi x'.\varphi x''\mathord{.}\,\mathord{.}\,\mathord{.}\,\mathord{.}\,\mathrm{d}y.\mathrm{d}x'.\mathrm{d}x''\mathord{.}\,\mathord{.}\,\mathord{.}\,\mathord{.}
\end{aligned}
\]
the integration extending from \(y=\eta\) to \(y=\eta'\) and over all values of the indeterminates \(x',x''\) etc. to which a real value of \(f(y,x',x''\mathord{.}\,\mathord{.}\,\mathord{.}\,\mathord{.})\) corresponds; or, what is the same, to the value of the integral
\[
\int y\varphi x.\varphi x'.\varphi x''\mathord{.}\,\mathord{.}\,\mathord{.}\,\mathord{.}\,\mathrm{d}x.\mathrm{d}x'.\mathrm{d}x''\mathord{.}\,\mathord{.}\,\mathord{.}\,\mathord{.}
\]
using in this integration for \(y\) its value expressed in terms of \(x,x',x''\) etc., and extending the integration over all values of these indeterminates to which a value of \(y\) lying between \(\eta\) and \(\eta'\) corresponds. From this we conclude that the integral \(\int y\psi y.\mathrm{d}y\) extended over all values of \(y\), from \(y=-\infty\) to \(y=+\infty\), is obtained from the integration of
\[
\int y\varphi x.\varphi x'.\varphi x''\mathord{.}\,\mathord{.}\,\mathord{.}\,\mathord{.}\,\mathrm{d}x.\mathrm{d}x'.\mathrm{d}x''\mathord{.}\,\mathord{.}\,\mathord{.}\,\mathord{.}
\]
\SourceLeaf{019}extended over all real values of \(x,x',x''\) etc., namely from \(x=-\infty\) to \(x=+\infty\), from \(x'=-\infty\) to \(x'=+\infty\) etc.

\clearpage
